% Options for packages loaded elsewhere
\PassOptionsToPackage{unicode}{hyperref}
\PassOptionsToPackage{hyphens}{url}
\documentclass[
  justified]{tufte-book}
\usepackage{xcolor}
\usepackage{amsmath,amssymb}
\setcounter{secnumdepth}{-\maxdimen} % remove section numbering
\usepackage{iftex}
\ifPDFTeX
  \usepackage[T1]{fontenc}
  \usepackage[utf8]{inputenc}
  \usepackage{textcomp} % provide euro and other symbols
\else % if luatex or xetex
  \usepackage{unicode-math} % this also loads fontspec
  \defaultfontfeatures{Scale=MatchLowercase}
  \defaultfontfeatures[\rmfamily]{Ligatures=TeX,Scale=1}
\fi
\usepackage{lmodern}
\ifPDFTeX\else
  % xetex/luatex font selection
\fi
% Use upquote if available, for straight quotes in verbatim environments
\IfFileExists{upquote.sty}{\usepackage{upquote}}{}
\IfFileExists{microtype.sty}{% use microtype if available
  \usepackage[]{microtype}
  \UseMicrotypeSet[protrusion]{basicmath} % disable protrusion for tt fonts
}{}
\makeatletter
\@ifundefined{KOMAClassName}{% if non-KOMA class
  \IfFileExists{parskip.sty}{%
    \usepackage{parskip}
  }{% else
    \setlength{\parindent}{0pt}
    \setlength{\parskip}{6pt plus 2pt minus 1pt}}
}{% if KOMA class
  \KOMAoptions{parskip=half}}
\makeatother
\usepackage{color}
\usepackage{fancyvrb}
\newcommand{\VerbBar}{|}
\newcommand{\VERB}{\Verb[commandchars=\\\{\}]}
\DefineVerbatimEnvironment{Highlighting}{Verbatim}{commandchars=\\\{\}}
% Add ',fontsize=\small' for more characters per line
\newenvironment{Shaded}{}{}
\newcommand{\AlertTok}[1]{\textcolor[rgb]{1.00,0.00,0.00}{\textbf{#1}}}
\newcommand{\AnnotationTok}[1]{\textcolor[rgb]{0.38,0.63,0.69}{\textbf{\textit{#1}}}}
\newcommand{\AttributeTok}[1]{\textcolor[rgb]{0.49,0.56,0.16}{#1}}
\newcommand{\BaseNTok}[1]{\textcolor[rgb]{0.25,0.63,0.44}{#1}}
\newcommand{\BuiltInTok}[1]{\textcolor[rgb]{0.00,0.50,0.00}{#1}}
\newcommand{\CharTok}[1]{\textcolor[rgb]{0.25,0.44,0.63}{#1}}
\newcommand{\CommentTok}[1]{\textcolor[rgb]{0.38,0.63,0.69}{\textit{#1}}}
\newcommand{\CommentVarTok}[1]{\textcolor[rgb]{0.38,0.63,0.69}{\textbf{\textit{#1}}}}
\newcommand{\ConstantTok}[1]{\textcolor[rgb]{0.53,0.00,0.00}{#1}}
\newcommand{\ControlFlowTok}[1]{\textcolor[rgb]{0.00,0.44,0.13}{\textbf{#1}}}
\newcommand{\DataTypeTok}[1]{\textcolor[rgb]{0.56,0.13,0.00}{#1}}
\newcommand{\DecValTok}[1]{\textcolor[rgb]{0.25,0.63,0.44}{#1}}
\newcommand{\DocumentationTok}[1]{\textcolor[rgb]{0.73,0.13,0.13}{\textit{#1}}}
\newcommand{\ErrorTok}[1]{\textcolor[rgb]{1.00,0.00,0.00}{\textbf{#1}}}
\newcommand{\ExtensionTok}[1]{#1}
\newcommand{\FloatTok}[1]{\textcolor[rgb]{0.25,0.63,0.44}{#1}}
\newcommand{\FunctionTok}[1]{\textcolor[rgb]{0.02,0.16,0.49}{#1}}
\newcommand{\ImportTok}[1]{\textcolor[rgb]{0.00,0.50,0.00}{\textbf{#1}}}
\newcommand{\InformationTok}[1]{\textcolor[rgb]{0.38,0.63,0.69}{\textbf{\textit{#1}}}}
\newcommand{\KeywordTok}[1]{\textcolor[rgb]{0.00,0.44,0.13}{\textbf{#1}}}
\newcommand{\NormalTok}[1]{#1}
\newcommand{\OperatorTok}[1]{\textcolor[rgb]{0.40,0.40,0.40}{#1}}
\newcommand{\OtherTok}[1]{\textcolor[rgb]{0.00,0.44,0.13}{#1}}
\newcommand{\PreprocessorTok}[1]{\textcolor[rgb]{0.74,0.48,0.00}{#1}}
\newcommand{\RegionMarkerTok}[1]{#1}
\newcommand{\SpecialCharTok}[1]{\textcolor[rgb]{0.25,0.44,0.63}{#1}}
\newcommand{\SpecialStringTok}[1]{\textcolor[rgb]{0.73,0.40,0.53}{#1}}
\newcommand{\StringTok}[1]{\textcolor[rgb]{0.25,0.44,0.63}{#1}}
\newcommand{\VariableTok}[1]{\textcolor[rgb]{0.10,0.09,0.49}{#1}}
\newcommand{\VerbatimStringTok}[1]{\textcolor[rgb]{0.25,0.44,0.63}{#1}}
\newcommand{\WarningTok}[1]{\textcolor[rgb]{0.38,0.63,0.69}{\textbf{\textit{#1}}}}
% definitions for citeproc citations
\NewDocumentCommand\citeproctext{}{}
\NewDocumentCommand\citeproc{mm}{%
  \begingroup\def\citeproctext{#2}\cite{#1}\endgroup}
\makeatletter
 % allow citations to break across lines
 \let\@cite@ofmt\@firstofone
 % avoid brackets around text for \cite:
 \def\@biblabel#1{}
 \def\@cite#1#2{{#1\if@tempswa , #2\fi}}
\makeatother
\newlength{\cslhangindent}
\setlength{\cslhangindent}{1.5em}
\newlength{\csllabelwidth}
\setlength{\csllabelwidth}{3em}
\newenvironment{CSLReferences}[2] % #1 hanging-indent, #2 entry-spacing
 {\begin{list}{}{%
  \setlength{\itemindent}{0pt}
  \setlength{\leftmargin}{0pt}
  \setlength{\parsep}{0pt}
  % turn on hanging indent if param 1 is 1
  \ifodd #1
   \setlength{\leftmargin}{\cslhangindent}
   \setlength{\itemindent}{-1\cslhangindent}
  \fi
  % set entry spacing
  \setlength{\itemsep}{#2\baselineskip}}}
 {\end{list}}
\usepackage{calc}
\newcommand{\CSLBlock}[1]{\hfill\break\parbox[t]{\linewidth}{\strut\ignorespaces#1\strut}}
\newcommand{\CSLLeftMargin}[1]{\parbox[t]{\csllabelwidth}{\strut#1\strut}}
\newcommand{\CSLRightInline}[1]{\parbox[t]{\linewidth - \csllabelwidth}{\strut#1\strut}}
\newcommand{\CSLIndent}[1]{\hspace{\cslhangindent}#1}
\setlength{\emergencystretch}{3em} % prevent overfull lines
\providecommand{\tightlist}{%
  \setlength{\itemsep}{0pt}\setlength{\parskip}{0pt}}
\makeatletter\newcommand{\subtitle}[1]{\gdef\@subtitle{#1}}\let\@subtitle\@empty\makeatother
\makeatletter\let\@oldmaketitle\maketitle\renewcommand{\maketitle}{\@oldmaketitle\ifx\@subtitle\@empty\else{\bigskip\noindent\Large\itshape\@subtitle\par}\fi}\makeatother
\usepackage{fvextra}
\DefineVerbatimEnvironment{Highlighting}{Verbatim}{breaklines,commandchars=\\\{\}}
\hypersetup{colorlinks=true, linkcolor=black, urlcolor=black, citecolor=black}
\usepackage{bookmark}
\IfFileExists{xurl.sty}{\usepackage{xurl}}{} % add URL line breaks if available
\urlstyle{same}
\hypersetup{
  pdftitle={Simple Ain't Stupid},
  pdfauthor={Tim Menzies},
  hidelinks,
  pdfcreator={LaTeX via pandoc}}

\title{Simple Ain't Stupid}
\usepackage{etoolbox}
\makeatletter
\providecommand{\subtitle}[1]{% add subtitle to \maketitle
  \apptocmd{\@title}{\par {\large #1 \par}}{}{}
}
\makeatother
\subtitle{200 Lessons in SE, AI, and Python from 2000 Lines of Code}
\author{Tim Menzies}
\date{2026 (draft)}

\begin{document}
\frontmatter
\maketitle

\mainmatter
\chapter*{Preface}\label{preface}
\addcontentsline{toc}{chapter}{Preface}

Much of what we do in software engineering and AI comes from a small
number of core ideas, wrapped in vast amounts of ceremony. This book
tries to show the core without the ceremony. So we set ourselves a task:
teach 200 heuristics, tips, tricks, and traps from scripting, SE, and
AI, using as little code as possible. Chasing that task produced EZR (an
incremental active learner whose conclusions are small decision trees),
about 400 lines of Python with no dependencies beyond the standard
library (\citeproc{ref-menzies26ezr}{Menzies and Srinivasan 2026}). We
do not say EZR is the only good toolkit. It is the one this hunt left
us, and it is small enough to read. Around it we hang 123 lessons. Many
appear more than once, seen from different heights, so the count of
sightings runs past 200.

In the age of large language models, it is said that software is free:
ask, and code appears. We agree, in part. Most software is now cheap.
Good software is very, very expensive. Good software is maintainable,
clear enough for a stranger to understand, and built so the next feature
has somewhere to go. Good software is not rushed out the door and left
to gather dirt. It is inspected with care, cleaned regularly, and
refactored, over and over, to be made simpler. That work is the
expensive part, and that work is what this book teaches.

Every chapter follows the same ritual: (a) set a seed; (b) run a small
program; (c) paste its output, untouched; (d) end in an assert. Nothing
you read here was typed from memory. The build tool re-runs every
example and fails if the output drifts.\footnote{If you find a
  transcript in this book that does not reproduce, that is a bug. Please
  report it.}

\section*{How to read this book}\label{how-to-read-this-book}
\addcontentsline{toc}{section}{How to read this book}

Read Chapter 1, then skip Chapter 2. That chapter places this book
against fifty years of prior work and it will mean more after you have
met the code. Come back to it late, or never.

Chapters 3 to 5 build the substrate: a little Python, a little maths,
then columns, tables, and distance. All later chapters stand on these
three. After that, each chapter is one application with a fun name and a
serious bracket (e.g.~``The Bouncer (anomaly detection)''). The
bracketed term is the index entry that joins that chapter to the
standard literature. Read the applications in any order.

You need Python 3.12 or later, curl, and make. You do not need pip. The
number of packages this book installs is 0 + 0 = 0. Example data comes
from MOOT (a public collection of 120+ multi-objective SE optimization
tasks)\footnote{\texttt{github.com/timm/moot}; fetched by
  \texttt{make\ data}.} and one table runs through most demos: auto93,
398 cars, where we want to minimize weight while maximizing acceleration
and miles per gallon.

\part{The Substrate}

\chapter{Simple Ain't Stupid}\label{simple-aint-stupid}

Much recent press says that developers no longer need to read code. The
creator of Node.js suggests the era of human-written code is ending.
Nvidia's CEO advises against learning to code at all
(\citeproc{ref-menzies26ezr}{Menzies and Srinivasan 2026}). The claim
behind the phrasing is always the same: AI is the new compiler, and
nobody reads what a compiler emits.

We disagree. This book is our evidence.

The evidence is 400 lines of Python that implement Naive Bayes, k-means
clustering, decision trees, anomaly detection, active learning, and
multi-objective optimization on one shared substrate. Tested on 120+
tasks from the MOOT repository, these tiny tools perform as well as or
better than SHAP, LIME, and the SMAC3 optimizer, while running 500 times
faster than SMAC3 and using orders of magnitude fewer labels
(\citeproc{ref-menzies26ezr}{Menzies and Srinivasan 2026}). That result
was found by reading: years of taking learners apart, noticing that
distant parts of different algorithms were the same part, and deleting
the copies. Reading code is not nostalgia. It is a research method, and
it works.

\section{The axiom}\label{the-axiom}

One economic fact organizes everything that follows. Tables of data are
cheap; \emph{labels} are dear. Anyone can list 10,000 configuration
options. Knowing which one compiles fastest costs a compile per option.
Hence the question this book asks, over and over: how few labels buy a
good row, plus an explanation of why it is good?

Our running example is auto93: 398 cars, 8 columns. Four columns are
cheap to observe (cylinders, engine volume, model year, origin). Three
are goals we must pay to know: minimize Lbs-, maximize Acc+ and Mpg+.
One (HpX) we ignore. That is 4 + 3 + 1 = 8. When we say a method is
good, we mean it finds near-best cars after paying for few labels.

\section{The ritual}\label{the-ritual}

\textbf{Every claim in this book can be re-run from a seed, and any
claim that cannot is not in this book.} Each demo (a) resets the random
stream; (b) runs; (c) prints; (d) asserts. The printed output is
captured by the build tool at build time, so what you read is what the
code did, on the day this book was compiled. Chapter 15 shows the
statistics that police the stronger claims: no ``X beats Y'' without
effect size and significance agreeing.

\section{The map}\label{the-map}

Part I (Chapters 3 to 5) builds the substrate: cells, columns, tables,
distance. Part II reads the world: six applications that predict,
detect, monitor, diagnose, triage, and explain. Part III changes the
world: eight applications that optimize, plan, repair, generate,
compress, negotiate, and learn on streams. Part IV earns trust:
certification by statistics, and the onboarding of an AI colleague.
Chapter 2 relates all this to prior work. Skip it for now.

One convention, used throughout. Where a lesson has a canonical name,
that name appears in bold at first use (e.g.~\textbf{SSOT}, single
source of truth). Where a lesson has no canonical name, none is
invented. The unnamed ones are the new ones.

\chapter{Before Us (Skip This on First
Read)}\label{before-us-skip-this-on-first-read}

This chapter has one job: to say what is old here, so that what is new
stands out honestly. It cites much and proves nothing. Come back after
Part II.

\section{Four crowded neighborhoods}\label{four-crowded-neighborhoods}

Code with commentary, on one system, has a patron saint: Lions'
commentary on UNIX, the complete Edition 6 source plus a reading of it,
bootlegged for decades and still used in teaching
(\citeproc{ref-lions77}{Lions 1996}). Spinellis later made code reading
a discipline of its own (\citeproc{ref-spinellis03}{Spinellis 2003}). We
take from Lions the core structural move: print the system once, then
tour it repeatedly at different altitudes.

Language books with practicals form the second neighborhood. Norvig's
\emph{PAIP} taught classic AI through small idiomatic Common Lisp
programs (\citeproc{ref-norvig92}{Norvig 1992}). Seibel's
\emph{Practical Common Lisp} taught a language by building a spam
filter, an MP3 database, and other things a reader might actually want
(\citeproc{ref-seibel05}{Seibel 2005}). Our Part II and Part III copy
Seibel's test directly: every chapter is named for a want, and the
machinery hides behind it.

Small-programs anthologies are the third. \emph{500 Lines or Less} walks
26 programs by 26 authors, each solving a canonical problem in under 500
lines (\citeproc{ref-brown16}{Brown and DiBernardo 2016}). Wilson's
\emph{Software Design by Example} does the same solo
(\citeproc{ref-wilson22}{Wilson 2024}). These books share our size
discipline. They do not share our substrate: their programs have nothing
in common, so their lessons cannot compose. Ours do, since every
application here reuses the same 400 lines.

Heuristics catalogs are the fourth: Hunt and Thomas
(\citeproc{ref-hunt99}{Hunt and Thomas 1999}), Glass's facts and
fallacies (\citeproc{ref-glass02}{Glass 2002}), Ousterhout
(\citeproc{ref-ousterhout18}{Ousterhout 2018}). These argue by anecdote
and citation. We argue by flag: when this book says ``always keep a
baseline'', that baseline is one command-line option away, and you can
run the ablation yourself.

Each neighborhood is crowded. Their intersection, as far as we can tell,
is empty. That intersection is this book.

\section{The grid and the ladder}\label{the-grid-and-the-ladder}

Buse and Zimmermann once organized software analytics as a 3-by-3 grid:
exploration, analysis, and experimentation crossed with past, present,
and future (\citeproc{ref-buse12}{Buse and Zimmermann 2012}). It was,
and is, a useful catalog of what managers ask for. Yet we read the grid
differently. Its three rows are the three rungs of Pearl's ladder of
causation: seeing, explaining, and imagining alternatives
(\citeproc{ref-pearl18}{Pearl and Mackenzie 2018}). Its three columns
are a clock. And a clock is a presentation dimension, not a reasoning
one. Chapters 8, 9, and 10 of this book make that point in code:
planning, diagnosis, and repair turn out to be one tree-difference
mechanism pointed at the future, the past, and the imperative. Time
multiplies interfaces. It does not multiply mechanisms. Hence, in this
book, streaming is a crosscutting section in many chapters rather than a
wing of the taxonomy.

\section{Tasks, not tenses}\label{tasks-not-tenses}

For the taxonomy of tasks we reach further back, to knowledge
engineering. Clancey observed that expert systems mostly perform
\emph{heuristic classification} (\citeproc{ref-clancey85}{Clancey
1985}), and the CommonKADS school cataloged the task types
(\citeproc{ref-schreiber00}{Schreiber et al. 2000}): analysis tasks
(predict, monitor, diagnose, classify, explain) and synthesis tasks
(plan, design, configure, repair). That catalog is finite. Part II
covers the analysis tasks. Part III covers the synthesis tasks, plus the
modern additions (compress, negotiate, stream). We know of no other
short book that runs the whole KADS catalog on one substrate.

There is one rung the old catalogs lack. Neither the grid nor the ladder
nor KADS has a place for \emph{certification}: establishing that the
seeing and the doing can be trusted by another person, another runtime,
or another kind of colleague. Buse and Zimmermann tuck significance
testing into a corner cell as a bullet point. We give it Part IV. The
2026 problem, we will argue, lives on that rung, and it costs about
sixty lines.

\section{What is actually new}\label{what-is-actually-new}

Firstly, the substrate: many seemingly different learners shown as one
grouping idea at different granularities, in runnable form. Secondly,
the receipts: heuristics with flags and seeds, not heuristics with
anecdotes. Thirdly, the last two chapters: a statistics gate and an
agent-onboarding process, both demonstrated from one repository's own
history. To our knowledge the third item exists in no prior book at all.
That said, it may date the fastest. We accept the trade.

\chapter{A Little Python}\label{a-little-python}

This chapter shows the Python that the rest of the book assumes. Not all
of Python. Just the dozen moves that let 400 lines carry twenty
algorithms. Readers fluent in the language should skim the transcripts
and move on.

\section{Cells, and asking
forgiveness}\label{cells-and-asking-forgiveness}

Csv cells arrive as strings. One function coerces each cell: ``23''
becomes an int, ``-1e2'' a float (csv cells can hide exponents),
``True'' and ``False'' become bools, ``?'' stays as our mark for
missing, and anything else stays text.

\begin{Shaded}
\begin{Highlighting}[]
\KeywordTok{def}\NormalTok{ thing(s):}
  \CommentTok{"Coerce one csv cell: int, float, bool, else string."}
  \ControlFlowTok{for}\NormalTok{ fn }\KeywordTok{in}\NormalTok{ (}\BuiltInTok{int}\NormalTok{, }\BuiltInTok{float}\NormalTok{):}
    \ControlFlowTok{try}\NormalTok{: }\ControlFlowTok{return}\NormalTok{ fn(s)}
    \ControlFlowTok{except} \PreprocessorTok{ValueError}\NormalTok{: }\ControlFlowTok{pass}
\NormalTok{  s }\OperatorTok{=}\NormalTok{ s.strip()}
  \ControlFlowTok{return}\NormalTok{ \{}\StringTok{"True"}\NormalTok{: }\VariableTok{True}\NormalTok{, }\StringTok{"False"}\NormalTok{: }\VariableTok{False}\NormalTok{\}.get(s, s)}
\end{Highlighting}
\end{Shaded}

Note the shape of the type test. We do not inspect the string to decide
if it is a number. We try the conversion and catch the failure. Python
calls this \textbf{EAFP} (easier to ask forgiveness than permission), as
opposed to \textbf{LBYL} (look before you leap). EAFP is shorter here
and, more importantly, it delegates the definition of ``looks like an
int'' to the one authority that matters, which is \texttt{int} itself.

\begin{verbatim}
$ python3 src/lib_eg.py thing
[23, 3.14, -100.0, True, False, '?', 'ab']
\end{verbatim}

\section{One struct, one settings
object}\label{one-struct-one-settings-object}

The whole book uses one generic record type: a class \texttt{o} whose
instances are dot-accessible bags of slots, and whose repr prints those
slots sorted.

\begin{Shaded}
\begin{Highlighting}[]
\KeywordTok{class}\NormalTok{ o:}
  \CommentTok{"Dot{-}access struct. Its repr prints the public slots."}
  \KeywordTok{def} \FunctionTok{\_\_init\_\_}\NormalTok{(i, }\OperatorTok{**}\NormalTok{d): i.\_\_dict\_\_.update(}\OperatorTok{**}\NormalTok{d)}
  \KeywordTok{def} \FunctionTok{\_\_repr\_\_}\NormalTok{(i):}
    \ControlFlowTok{return} \StringTok{"\{"} \OperatorTok{+} \StringTok{" "}\NormalTok{.join(}\StringTok{":}\SpecialCharTok{\%s}\StringTok{ }\SpecialCharTok{\%s}\StringTok{"} \OperatorTok{\%}\NormalTok{ (k, v)}
      \ControlFlowTok{for}\NormalTok{ k, v }\KeywordTok{in} \BuiltInTok{sorted}\NormalTok{(i.\_\_dict\_\_.items())}
      \ControlFlowTok{if}\NormalTok{ k[}\DecValTok{0}\NormalTok{] }\OperatorTok{!=} \StringTok{"\_"}\NormalTok{) }\OperatorTok{+} \StringTok{"\}"}
\end{Highlighting}
\end{Shaded}

Two Python facts do the work: (a) every object carries a
\texttt{\_\_dict\_\_} of its slots, so one line of reflection prints
anything; (b) \texttt{\_\_repr\_\_} is a protocol, so every object in
this book can appear in a transcript. Also, all settings live in a
single instance called \texttt{the}, defined in about.py and nowhere
else. That file is this book's first sighting of \textbf{SSOT} (single
source of truth): one place to look, one place to change.

\begin{verbatim}
$ python3 src/lib_eg.py the
{:few 128 :file data/auto93.csv :p 2 :seed 1234567891}
\end{verbatim}

The command line can override any knob, because the flag parser walks
\texttt{vars(the)} and matches names. Watch the same test, run with
\texttt{-\/-p=1}:

\begin{verbatim}
$ python3 src/lib_eg.py the --p=1
{:few 128 :file data/auto93.csv :p 1 :seed 1234567891}
\end{verbatim}

No parser was written for that flag. The settings object is the parser's
schema, by reflection. Hence a new knob costs one line, in one file.

\section{Streams, sorts, and
transposes}\label{streams-sorts-and-transposes}

Files are read by a generator, so no table is ever loaded twice or held
whole when a stream will do:

\begin{Shaded}
\begin{Highlighting}[]
\KeywordTok{def}\NormalTok{ csv(}\BuiltInTok{file}\NormalTok{):}
  \CommentTok{"Stream a csv file, one row of coerced cells at a time."}
  \ControlFlowTok{with} \BuiltInTok{open}\NormalTok{(}\BuiltInTok{file}\NormalTok{) }\ImportTok{as}\NormalTok{ f:}
    \ControlFlowTok{for}\NormalTok{ line }\KeywordTok{in}\NormalTok{ f:}
\NormalTok{      line }\OperatorTok{=}\NormalTok{ line.split(}\StringTok{"\%"}\NormalTok{)[}\DecValTok{0}\NormalTok{].strip()}
      \ControlFlowTok{if}\NormalTok{ line:}
        \ControlFlowTok{yield}\NormalTok{ [thing(s) }\ControlFlowTok{for}\NormalTok{ s }\KeywordTok{in}\NormalTok{ line.split(}\StringTok{","}\NormalTok{)]}
\end{Highlighting}
\end{Shaded}

The \texttt{yield} makes \texttt{csv} lazy: callers pull one row at a
time. Later chapters lean hard on this laziness (e.g.~the streaming
chapter treats an unbounded source exactly like a file). Three more
idioms appear on nearly every page of this book, so we test them once
here: sorting by a computed key, never in place; transposing a table
with \texttt{zip(*rows)}; and list comprehensions as the default loop.

\begin{verbatim}
$ python3 src/lib_eg.py idioms
[[3, 30], [2, 20], [1, 10]]
[(1, 3, 2), (10, 30, 20)]
\end{verbatim}

\section{Dispatch by name}\label{dispatch-by-name}

Every code file in this book ships with a matching \texttt{\_eg} file of
demos. Each demo runs by its bare name from the command line. The trick
is four lines of reflection. The caller hands over its
\texttt{globals()} as \texttt{g}, a plain dict mapping names to the
things they name. On any dict, \texttt{g.get(key,\ default)} returns
\texttt{g{[}key{]}} if the key is present, else the default. Hence: look
up \texttt{"test\_"\ +\ word}, falling back to a function that just
complains.

\begin{Shaded}
\begin{Highlighting}[]
\KeywordTok{def}\NormalTok{ main(g):}
  \CommentTok{"For each bare command{-}line word w, run test\_w, seeded."}
\NormalTok{  cli(the.\_\_dict\_\_)}
\NormalTok{  todo }\OperatorTok{=}\NormalTok{ [s }\ControlFlowTok{for}\NormalTok{ s }\KeywordTok{in}\NormalTok{ sys.argv[}\DecValTok{1}\NormalTok{:]}
          \ControlFlowTok{if} \KeywordTok{not}\NormalTok{ s.startswith(}\StringTok{"{-}"}\NormalTok{)] }\KeywordTok{or}\NormalTok{ [}\StringTok{"all"}\NormalTok{]}
  \ControlFlowTok{for}\NormalTok{ word }\KeywordTok{in}\NormalTok{ todo:}
\NormalTok{    random.seed(the.seed)}
\NormalTok{    g.get(}\StringTok{"test\_"} \OperatorTok{+}\NormalTok{ word,}
          \KeywordTok{lambda}\NormalTok{: }\BuiltInTok{print}\NormalTok{(}\StringTok{"?"}\NormalTok{, word, }\StringTok{"(no such test)"}\NormalTok{))()}
\end{Highlighting}
\end{Shaded}

Note that the seed is reset before every test. Hence any demo reproduces
in isolation, in any order, on any machine. This one habit does more for
reproducibility than any tool we know, and it costs one line.

\section{Lessons sighted}\label{lessons-sighted}

\textbf{EAFP} and \textbf{LBYL}; \textbf{SSOT}; duck typing (previewed;
defined properly in Chapter 5); generators and laziness; reflection via
\texttt{vars} and \texttt{globals}; sort-by-key; \texttt{zip(*rows)};
the seeded-demo ritual. Just to repeat a point made above: none of this
is advanced Python. That is the point. Relentless basics, then one or
two sharp tricks.

\chapter{A Little Maths}\label{a-little-maths}

Six pieces of maths run the whole book: two summaries, a rescaling, a
distance, a random stream, and a trust test. None needs more than
high-school algebra. Each gets a demo here and a citation for readers
who want the long form. But first, some words.

\section{Four kinds of column, two
survive}\label{four-kinds-of-column-two-survive}

Statistics sorts columns into four scales, by what you may do to their
values: nominal (names; test equality, nothing else), ordinal (ranks;
also sort), interval (also subtract), and ratio (also divide)
(\citeproc{ref-stevens46}{Stevens 1946}). The acronym is \textbf{NOIR}.
This book keeps the two ends. Names become \texttt{Sym} and numbers
become \texttt{Num}; the middle two get treated as one or the other.

\begin{Shaded}
\begin{Highlighting}[]
\KeywordTok{def}\NormalTok{ Sym(at}\OperatorTok{=}\DecValTok{0}\NormalTok{, name}\OperatorTok{=}\StringTok{" "}\NormalTok{):}
  \CommentTok{"Summary of a symbolic column."}
  \ControlFlowTok{return}\NormalTok{ o(it}\OperatorTok{=}\NormalTok{Sym, at}\OperatorTok{=}\NormalTok{at, name}\OperatorTok{=}\NormalTok{name, n}\OperatorTok{=}\DecValTok{0}\NormalTok{, has}\OperatorTok{=}\NormalTok{\{\})}
\end{Highlighting}
\end{Shaded}

\begin{Shaded}
\begin{Highlighting}[]
\KeywordTok{def}\NormalTok{ Num(at}\OperatorTok{=}\DecValTok{0}\NormalTok{, name}\OperatorTok{=}\StringTok{" "}\NormalTok{):}
  \CommentTok{"Summary of a numeric column."}
  \ControlFlowTok{return}\NormalTok{ o(it}\OperatorTok{=}\NormalTok{Num, at}\OperatorTok{=}\NormalTok{at, name}\OperatorTok{=}\NormalTok{name, n}\OperatorTok{=}\DecValTok{0}\NormalTok{, mu}\OperatorTok{=}\DecValTok{0}\NormalTok{, m2}\OperatorTok{=}\DecValTok{0}\NormalTok{,}
\NormalTok{           heaven}\OperatorTok{=}\DecValTok{0} \ControlFlowTok{if}\NormalTok{ name.endswith(}\StringTok{"{-}"}\NormalTok{) }\ControlFlowTok{else} \DecValTok{1}\NormalTok{)}
\end{Highlighting}
\end{Shaded}

Every column summary answers the same two questions. Where is the
middle? That is central tendency: the mean mu for \texttt{Num}, the mode
(the commonest symbol) for \texttt{Sym}. How far do values stray from
that middle? That is dispersion: the standard deviation for
\texttt{Num}, the entropy for \texttt{Sym}. The next two sections build
those four answers, one value at a time.

Two more words before we start. A \textbf{pdf} (probability density
function) scores how likely each value is; its familiar picture is the
bell curve. A \textbf{cdf} (cumulative distribution function) reports
what fraction of the data sits at or below v. Whatever the units, a cdf
runs 0..1. Remember that; it is about to earn its keep.

\section{Center and spread, one value at a
time}\label{center-and-spread-one-value-at-a-time}

For numbers we track the mean and the standard deviation. The textbook
formula for variance wants two passes and a stored list. Welford's
update (\citeproc{ref-welford62}{Welford 1962}) wants neither. Keep
three slots (n, mu, m2). On each new value v, let d = v - mu. Then mu
moves by d/n, and m2 grows by d times the \emph{new} gap (v - mu). The
standard deviation is (m2 / (n - 1)) raised to 0.5, on demand.

\begin{Shaded}
\begin{Highlighting}[]
\KeywordTok{def}\NormalTok{ welford(num, v, inc}\OperatorTok{=}\DecValTok{1}\NormalTok{):}
  \CommentTok{"One{-}pass mean and spread update; inc={-}1 forgets."}
\NormalTok{  num.n }\OperatorTok{+=}\NormalTok{ inc}
  \ControlFlowTok{if}\NormalTok{ inc }\OperatorTok{\textless{}} \DecValTok{0} \KeywordTok{and}\NormalTok{ num.n }\OperatorTok{\textless{}} \DecValTok{2}\NormalTok{:}
\NormalTok{    num.mu }\OperatorTok{=}\NormalTok{ num.m2 }\OperatorTok{=}\NormalTok{ num.n }\OperatorTok{=} \DecValTok{0}
  \ControlFlowTok{else}\NormalTok{:}
\NormalTok{    d       }\OperatorTok{=}\NormalTok{ v }\OperatorTok{{-}}\NormalTok{ num.mu}
\NormalTok{    num.mu }\OperatorTok{+=}\NormalTok{ inc }\OperatorTok{*}\NormalTok{ d }\OperatorTok{/}\NormalTok{ num.n}
\NormalTok{    num.m2 }\OperatorTok{+=}\NormalTok{ inc }\OperatorTok{*}\NormalTok{ d }\OperatorTok{*}\NormalTok{ (v }\OperatorTok{{-}}\NormalTok{ num.mu)}
\end{Highlighting}
\end{Shaded}

One pass. No stored raws. Constant memory. Symbols are easier. Count
them:

\begin{Shaded}
\begin{Highlighting}[]
\KeywordTok{def}\NormalTok{ count(sym, v, inc}\OperatorTok{=}\DecValTok{1}\NormalTok{):}
  \CommentTok{"Update symbol counts; inc={-}1 downdates them."}
\NormalTok{  sym.n }\OperatorTok{+=}\NormalTok{ inc}
\NormalTok{  sym.has[v] }\OperatorTok{=}\NormalTok{ inc }\OperatorTok{+}\NormalTok{ sym.has.get(v, }\DecValTok{0}\NormalTok{)}
\end{Highlighting}
\end{Shaded}

Now the public face. \texttt{add} reads the column's \texttt{it} tag,
skips the ``?'' that marks a missing value, and hands the rest to the
right updater. All the cleverness lives in the parts. The whole is one
line of dispatch:

\begin{Shaded}
\begin{Highlighting}[]
\KeywordTok{def}\NormalTok{ add(col, v, inc}\OperatorTok{=}\DecValTok{1}\NormalTok{):}
  \CommentTok{"Fold v into col. Returns v."}
  \ControlFlowTok{if}\NormalTok{ v }\OperatorTok{!=} \StringTok{"?"}\NormalTok{:}
\NormalTok{    (count }\ControlFlowTok{if}\NormalTok{ col.it }\KeywordTok{is}\NormalTok{ Sym }\ControlFlowTok{else}\NormalTok{ welford)(col, v, inc)}
  \ControlFlowTok{return}\NormalTok{ v}
\end{Highlighting}
\end{Shaded}

Also, look at the \texttt{inc} argument in both updaters: the same code
run with inc = -1 makes a summary \emph{forget}. That trick earns a
name:

\begin{Shaded}
\begin{Highlighting}[]
\KeywordTok{def}\NormalTok{ sub(col, v):}
  \CommentTok{"The undo of add: col forgets v. Returns v."}
  \ControlFlowTok{return}\NormalTok{ add(col, v, }\OperatorTok{{-}}\DecValTok{1}\NormalTok{)}
\end{Highlighting}
\end{Shaded}

Since add and sub both return v, they compose. The idiom add(i, sub(j,
v)) pops a value from summary j and pushes it into summary i, so a value
walks between summaries in one line. Chapter 13 builds sliding windows
from nothing else. Here is the round trip, exact to nine decimals, then
the walk:

\begin{verbatim}
$ python3 src/lib_eg.py unadd
mu 9.948 sd 2.215 (before: 9.948 2.215)
\end{verbatim}

\begin{verbatim}
$ python3 src/lib_eg.py move
i: n 75 mu 6.236 sd 2.322
\end{verbatim}

\section{Entropy is just counting}\label{entropy-is-just-counting}

For a bag of symbols with counts k out of n, the entropy is the sum over
symbols of -(k/n) log2 (k/n). It measures surprise: how many yes-or-no
questions, on average, to name the next symbol. Take the bag
``aaaabbc'', so n = 7 and the counts are 4, 2, 1. The three terms are
-(4/7) log2 (4/7) = 0.461, -(2/7) log2 (2/7) = 0.516, and -(1/7) log2
(1/7) = 0.401. Their sum is 0.461 + 0.516 + 0.401 = 1.379
bits.\footnote{Purists will note we never round away the working. House
  rule.} In code, the two middles share one roof, and so do the two
spreads:

\begin{Shaded}
\begin{Highlighting}[]
\KeywordTok{def}\NormalTok{ mid(col):}
  \CommentTok{"Central tendency: mean (Num) or mode (Sym)."}
  \ControlFlowTok{return}\NormalTok{ col.mu }\ControlFlowTok{if}\NormalTok{ col.it }\KeywordTok{is}\NormalTok{ Num }\ControlFlowTok{else} \OperatorTok{\textbackslash{}}
         \BuiltInTok{max}\NormalTok{(col.has, key}\OperatorTok{=}\NormalTok{col.has.get)}
\end{Highlighting}
\end{Shaded}

\begin{Shaded}
\begin{Highlighting}[]
\KeywordTok{def}\NormalTok{ var(col):}
  \CommentTok{"Dispersion: standard deviation (Num) or entropy (Sym)."}
  \ControlFlowTok{if}\NormalTok{ col.it }\KeywordTok{is}\NormalTok{ Sym:}
    \ControlFlowTok{return} \OperatorTok{{-}}\BuiltInTok{sum}\NormalTok{(k}\OperatorTok{/}\NormalTok{col.n }\OperatorTok{*}\NormalTok{ math.log(k}\OperatorTok{/}\NormalTok{col.n, }\DecValTok{2}\NormalTok{)}
                \ControlFlowTok{for}\NormalTok{ k }\KeywordTok{in}\NormalTok{ col.has.values() }\ControlFlowTok{if}\NormalTok{ k }\OperatorTok{\textgreater{}} \DecValTok{0}\NormalTok{)}
  \ControlFlowTok{return} \DecValTok{0} \ControlFlowTok{if}\NormalTok{ col.n }\OperatorTok{\textless{}} \DecValTok{2} \ControlFlowTok{else}\NormalTok{ (}\BuiltInTok{max}\NormalTok{(col.m2, }\DecValTok{0}\NormalTok{)}
                              \OperatorTok{/}\NormalTok{ (col.n }\OperatorTok{{-}} \DecValTok{1}\NormalTok{)) }\OperatorTok{**} \FloatTok{0.5}
\end{Highlighting}
\end{Shaded}

The demo checks the entropy arithmetic above, then checks Welford
against 10,000 samples from a unit gaussian:

\begin{verbatim}
$ python3 src/lib_eg.py cols
sym mid a ent 1.379
num mu 0.007 sd 1.002
\end{verbatim}

\section{Nothing is comparable until it is
0..1}\label{nothing-is-comparable-until-it-is-0..1}

In auto93, weight runs in the thousands of pounds while acceleration
runs in the tens of seconds. Any distance computed on raw values would
be a weight measure wearing a distance costume. Hence, before comparing,
we normalize using the normal: map each number to its column's cdf, the
fraction of the bell curve at or below v, computed from mu and sd via
the error function.\footnote{An older habit rescales by the seen range,
  (v - lo) / (hi - lo). That wants two more slots and its own epsilon,
  and one wild outlier squashes everyone else into a corner. The cdf
  runs on what Welford already keeps.} Whatever the units, the result
runs 0..1.

\begin{Shaded}
\begin{Highlighting}[]
\KeywordTok{def}\NormalTok{ norm(col, v):}
  \CommentTok{"Map v to 0..1: the gaussian cdf of v (Nums only)."}
  \ControlFlowTok{if}\NormalTok{ v }\OperatorTok{==} \StringTok{"?"} \KeywordTok{or}\NormalTok{ col.it }\KeywordTok{is}\NormalTok{ Sym: }\ControlFlowTok{return}\NormalTok{ v}
  \ControlFlowTok{return} \FloatTok{0.5} \OperatorTok{*}\NormalTok{ (}\DecValTok{1} \OperatorTok{+}\NormalTok{ math.erf(}
\NormalTok{    (v }\OperatorTok{{-}}\NormalTok{ col.mu) }\OperatorTok{/}\NormalTok{ (var(col) }\OperatorTok{*} \DecValTok{2} \OperatorTok{**} \FloatTok{0.5} \OperatorTok{+}\NormalTok{ TINY)))}
\end{Highlighting}
\end{Shaded}

A tiny epsilon guards the degenerate column whose sd is zero. Numerical
hygiene of that kind (a TINY here, a max(m2, 0) there) is not fussiness.
It is where AI code goes to die, quietly.

\section{One distance, three
classics}\label{one-distance-three-classics}

With every gap scaled to 0..1, we aggregate gaps by the Minkowski
formula: the p-th root of the mean of the p-th powers
(\citeproc{ref-menzies26ezr}{Menzies and Srinivasan 2026}). One knob,
three famous distances: p = 1 is city-block, p = 2 is Euclidean, and
large p approaches Chebyshev (the max gap wins). For two gaps of 0.3 and
0.4 at p = 2, that is ((0.09 + 0.16) / 2) \^{} 0.5 = (0.125) \^{} 0.5 =
0.354. You will meet the code in Chapter 5, written out longhand inside
the two distance functions. Distance sits in every inner loop this book
runs, so those functions stay flat: one call to norm per gap, and
nothing else.

\section{Random streams are lab
equipment}\label{random-streams-are-lab-equipment}

Every stochastic demo in this book starts by seeding the random stream,
so every transcript reproduces. Treat the generator as lab equipment:
calibrated, logged, reset between experiments. When we later port code
across languages, we will go further and use Lehmer's portable generator
with multiplier 16807, which is 7 * 7 * 7 * 7 * 7 = 16807, i.e.~7\^{}5
(\citeproc{ref-park88}{Park and Miller 1988}). Same seed, same stream,
in every language. Five printed draws then certify a port before any
learner runs.

\section{The trust test}\label{the-trust-test}

Later chapters keep saying ``X beats Y''. Such talk is cheap, so we tax
it. Two lists of results are called the same unless three tests all
agree they differ: Cohen's rule (means closer than 0.35 pooled standard
deviations are the same) (\citeproc{ref-cohen88}{Cohen 1988}), Cliff's
delta (rank overlap above the 0.197 threshold is the same)
(\citeproc{ref-cliff93}{Cliff 1993}; \citeproc{ref-hess04}{Hess and
Kromrey 2004}), and Kolmogorov-Smirnov (cdf gap under the 5\% critical
value 1.36 sqrt((n + m) / nm) is the same)
(\citeproc{ref-massey51}{Massey 1951}). Watch the conjunction work on a
gaussian nudged by ever-larger shifts. A shift of 0.1 standard
deviations passes as same. From 0.3 on, the tests call it different:

\begin{verbatim}
$ python3 src/lib_eg.py stats
shift  same   cliffs cohen
 +0.0  True   0.00   True
 +0.1  True   0.11   True
 +0.3  False  0.21   True
 +0.5  False  0.34   False
 +1.0  False  0.56   False
 +2.0  False  0.82   False
\end{verbatim}

Note the humility this buys. Any single test can be gamed, and p-values
alone say nothing about effect size. Demanding agreement from an
effect-size test and a distribution test before crying ``difference''
makes our later victory claims conservative. When this book says
``beats'', it has cleared all three bars.

\section{Lessons sighted}\label{lessons-sighted-1}

\textbf{NOIR} and the two surviving scales; Welford's one-pass update;
reversible summaries (every add has a sub); entropy as counted surprise;
cdf normalization and epsilon hygiene; the Minkowski family; seeds as
lab equipment (with \textbf{PRNG} portability via 7\^{}5); and the
conjunctive trust test. That last one is this book's referee. It sits in
the substrate, before any learner, on purpose.

\chapter{The Substrate}\label{the-substrate}

This chapter assembles the pieces: a settings file, two column types, a
table, and two distances. Everything in Parts II to IV is a short
function over what this chapter builds. Total cost so far, including the
test rig: under 250 lines.

\section{about.py, in full}\label{about.py-in-full}

\begin{Shaded}
\begin{Highlighting}[]
\CommentTok{\#!/usr/bin/env python3 {-}B}
\CommentTok{"""}
\CommentTok{about.py: every knob, one place. Change a value here, or}
\CommentTok{override it on the command line (e.g. {-}{-}p 1), and all}
\CommentTok{downstream code obeys. No other file defines a setting.}
\CommentTok{"""}

\KeywordTok{class}\NormalTok{ o:}
  \CommentTok{"Dot{-}access struct. Its repr prints the public slots."}
  \KeywordTok{def} \FunctionTok{\_\_init\_\_}\NormalTok{(i, }\OperatorTok{**}\NormalTok{d): i.\_\_dict\_\_.update(}\OperatorTok{**}\NormalTok{d)}
  \KeywordTok{def} \FunctionTok{\_\_repr\_\_}\NormalTok{(i):}
    \ControlFlowTok{return} \StringTok{"\{"} \OperatorTok{+} \StringTok{" "}\NormalTok{.join(}\StringTok{":}\SpecialCharTok{\%s}\StringTok{ }\SpecialCharTok{\%s}\StringTok{"} \OperatorTok{\%}\NormalTok{ (k, v)}
      \ControlFlowTok{for}\NormalTok{ k, v }\KeywordTok{in} \BuiltInTok{sorted}\NormalTok{(i.\_\_dict\_\_.items())}
      \ControlFlowTok{if}\NormalTok{ k[}\DecValTok{0}\NormalTok{] }\OperatorTok{!=} \StringTok{"\_"}\NormalTok{) }\OperatorTok{+} \StringTok{"\}"}

\NormalTok{the }\OperatorTok{=}\NormalTok{ o(}
\NormalTok{  seed }\OperatorTok{=} \DecValTok{1234567891}\NormalTok{,        }\CommentTok{\# every random stream starts here}
\NormalTok{  p    }\OperatorTok{=} \DecValTok{2}\NormalTok{,                 }\CommentTok{\# minkowski coefficient}
\NormalTok{  few  }\OperatorTok{=} \DecValTok{128}\NormalTok{,               }\CommentTok{\# sample size for cheap guesses}
  \BuiltInTok{file} \OperatorTok{=} \StringTok{"data/auto93.csv"}\NormalTok{) }\CommentTok{\# default table (via MOOT)}
\end{Highlighting}
\end{Shaded}

That is the entire configuration system. To say that another way: every
experiment in this book is fully described by (a) this file, (b) any
--key=val overrides, and (c) one seed. When Chapter 16 hands this
repository to an AI colleague, that property will matter a great deal.

\section{Columns know their jobs}\label{columns-know-their-jobs}

The first csv row names the columns, and the names carry the schema. An
uppercase first letter means numeric. A trailing ``+'' or ``-'' marks a
goal to maximize or minimize. A trailing ``X'' means ignore. Everything
else is an observable x column. This is \textbf{CoC} (convention over
configuration): the data describes itself, and no schema file exists to
drift out of date. It is also schema \emph{on read}: types come from
data, not declarations.

\begin{Shaded}
\begin{Highlighting}[]
\KeywordTok{def}\NormalTok{ Tbl(src):}
  \CommentTok{"First row names the columns; the rest are data."}
\NormalTok{  src   }\OperatorTok{=} \BuiltInTok{iter}\NormalTok{(src)}
\NormalTok{  names }\OperatorTok{=} \BuiltInTok{next}\NormalTok{(src)}
\NormalTok{  cols  }\OperatorTok{=}\NormalTok{ [(Num }\ControlFlowTok{if}\NormalTok{ s[}\DecValTok{0}\NormalTok{].isupper() }\ControlFlowTok{else}\NormalTok{ Sym)(at, s)}
           \ControlFlowTok{for}\NormalTok{ at, s }\KeywordTok{in} \BuiltInTok{enumerate}\NormalTok{(names)]}
\NormalTok{  tbl   }\OperatorTok{=}\NormalTok{ o(it}\OperatorTok{=}\NormalTok{Tbl, names}\OperatorTok{=}\NormalTok{names, cols}\OperatorTok{=}\NormalTok{cols, rows}\OperatorTok{=}\NormalTok{[],}
\NormalTok{            x}\OperatorTok{=}\NormalTok{[c.at }\ControlFlowTok{for}\NormalTok{ c }\KeywordTok{in}\NormalTok{ cols}
               \ControlFlowTok{if}\NormalTok{ c.name[}\OperatorTok{{-}}\DecValTok{1}\NormalTok{] }\KeywordTok{not} \KeywordTok{in} \StringTok{"X+{-}"}\NormalTok{],}
\NormalTok{            y}\OperatorTok{=}\NormalTok{[c.at }\ControlFlowTok{for}\NormalTok{ c }\KeywordTok{in}\NormalTok{ cols }\ControlFlowTok{if}\NormalTok{ c.name[}\OperatorTok{{-}}\DecValTok{1}\NormalTok{] }\KeywordTok{in} \StringTok{"+{-}"}\NormalTok{])}
  \ControlFlowTok{for}\NormalTok{ row }\KeywordTok{in}\NormalTok{ src: addRow(tbl, row)}
  \ControlFlowTok{return}\NormalTok{ tbl}
\end{Highlighting}
\end{Shaded}

Two details deserve a look. Firstly, \texttt{Num} and \texttt{Sym} are
plain functions returning \texttt{o} structs tagged with an \texttt{it}
slot, and one \texttt{add} serves both. That is duck typing doing the
work inheritance is usually hired for: two types, one protocol (add,
mid, var), no class hierarchy. Secondly, \texttt{addRow} folds a row
into every column summary incrementally. The table never recomputes. It
only ever updates.

\begin{verbatim}
$ python3 src/lib_eg.py tbl
rows 398 |x| 4 |y| 3
Mpg+ mu 23.84 sd 8.34
\end{verbatim}

Read that transcript against the axiom of Chapter 1. The table has 398
rows, and 4 of its 8 columns are cheap x values while 3 are dear y
goals. Everything downstream respects that boundary: x is free to read,
y is metered.

\texttt{clone} deserves its one line of fame. Given any table and any
subset of rows, it re-learns a fresh summary over just those rows, using
the same header. Clusters, tree leaves, and sliding windows are all,
underneath, clones.

\begin{Shaded}
\begin{Highlighting}[]
\KeywordTok{def}\NormalTok{ clone(tbl, rows}\OperatorTok{=}\NormalTok{[]):}
  \CommentTok{"A fresh table with the same header; optionally refill."}
  \ControlFlowTok{return}\NormalTok{ Tbl([tbl.names] }\OperatorTok{+}\NormalTok{ rows)}
\end{Highlighting}
\end{Shaded}

\section{Distance to heaven}\label{distance-to-heaven}

Now the two workhorses. \texttt{distx} measures how \emph{different} two
rows are, reading only x columns, normalizing each gap to 0..1 and
handling ``?'' by assuming the worst case (an unknown value is scored as
far away as it could be, which keeps missingness from faking
similarity). \texttt{disty} measures how \emph{good} one row is: the
distance from its goal values to heaven, where heaven is 0 for a
minimize column and 1 for a maximize column. Zero is best.

\begin{Shaded}
\begin{Highlighting}[]
\KeywordTok{def}\NormalTok{ disty(tbl, row):}
  \CommentTok{"Distance from a row\textquotesingle{}s goals to heaven; 0=best, 1=worst."}
\NormalTok{  d }\OperatorTok{=}\NormalTok{ n }\OperatorTok{=} \DecValTok{0}
  \ControlFlowTok{for}\NormalTok{ col }\KeywordTok{in}\NormalTok{ (tbl.cols[a] }\ControlFlowTok{for}\NormalTok{ a }\KeywordTok{in}\NormalTok{ tbl.y):}
\NormalTok{    v }\OperatorTok{=}\NormalTok{ row[col.at]}
    \ControlFlowTok{if}\NormalTok{ v }\OperatorTok{!=} \StringTok{"?"}\NormalTok{:}
\NormalTok{      d }\OperatorTok{+=} \BuiltInTok{abs}\NormalTok{(norm(col, v) }\OperatorTok{{-}}\NormalTok{ col.heaven) }\OperatorTok{**}\NormalTok{ the.p}
\NormalTok{      n }\OperatorTok{+=} \DecValTok{1}
  \ControlFlowTok{return}\NormalTok{ (d }\OperatorTok{/}\NormalTok{ (n }\OperatorTok{+}\NormalTok{ TINY)) }\OperatorTok{**}\NormalTok{ (}\DecValTok{1} \OperatorTok{/}\NormalTok{ the.p)}
\end{Highlighting}
\end{Shaded}

This paragraph is the definition of \textbf{distance to heaven}:
collapse many objectives to one scalar by measuring each goal's
normalized gap to its ideal, then aggregating with Minkowski. No weights
to tune. All later chapters cite this paragraph rather than redefine it.
Note also what \texttt{disty} reads: y columns only. Hence counting
calls to \texttt{disty} counts the labels we bought, and every budget
claim in this book is auditable from that one chokepoint.

Sort all 398 cars by \texttt{disty} and print the five best, then the
five worst:

\begin{verbatim}
$ python3 src/lib_eg.py dist
Clndrs  Volume  HpX  Model  origin  Lbs-  Acc+  Mpg+  disty
     4      90   48     78       2  1985  21.5    40  0.073
     4      85   65     81       3  1975  19.4    40  0.085
     4      90   48     80       2  2085  21.7    40  0.087
     4      97   52     82       2  2130  24.6    40  0.094
     4      79   58     77       2  1825  18.6    40  0.095

     8     440  215     73       1  4735    11    10  0.961
     8     455  225     70       1  4425    10    10  0.962
     8     440  215     70       1  4312   8.5    10  0.964
     8     454  220     70       1  4354     9    10  0.964
     8     455  225     73       1  4951    11    10  0.965
\end{verbatim}

Notice the shape. Light, late, high-mpg cars float to the top. Big old
guzzlers sink. No learner has run yet. This is just geometry over one
table, and already it ranks a car lot. The whole of Part III is a set of
strategies for reaching those top rows while paying for as few
\texttt{disty} calls as possible.

\section{What the substrate buys}\label{what-the-substrate-buys}

Here ends Part I. We close with the claim the next eleven chapters must
now earn. Clustering will be grouping by distx. Nearest-neighbor
prediction will be distx with no fit step. Bayes will be grouping with
labels. Trees will be grouping, recursively. Anomalies will be rows far
from their group. Each arrives in roughly twenty lines, because the hard
parts (summaries, normalization, missing values, distance, statistics)
already live here, written once. The good news is that you have now read
the hard parts. The bad news is two-fold: the substrate is dull, and we
made you read it anyway. It was that or repeat it eleven times.

\section{Lessons sighted}\label{lessons-sighted-2}

\textbf{SSOT} again (about.py, in full, as the experiment record);
\textbf{CoC} and schema-on-read via header suffixes; duck typing, now
defined; incremental summaries and clone; worst-case handling of missing
values; \textbf{distance to heaven}, defined above; labels metered at
one chokepoint. Onward to Part II, where the Fortune Teller takes the
first appointment.

\protect\phantomsection\label{refs}
\begin{CSLReferences}{1}{1}
\bibitem[\citeproctext]{ref-brown16}
Brown, Amy, and Michael DiBernardo, eds. 2016. \emph{500 Lines or Less:
Experienced Programmers Solve Interesting Problems}. Lulu.

\bibitem[\citeproctext]{ref-buse12}
Buse, Raymond P. L., and Thomas Zimmermann. 2012. {``Information Needs
for Software Development Analytics.''} \emph{Proc. ICSE}, 987--96.

\bibitem[\citeproctext]{ref-clancey85}
Clancey, William J. 1985. {``Heuristic Classification.''}
\emph{Artificial Intelligence} 27 (3): 289--350.

\bibitem[\citeproctext]{ref-cliff93}
Cliff, Norman. 1993. {``Dominance Statistics: Ordinal Analyses to Answer
Ordinal Questions.''} \emph{Psychological Bulletin} 114 (3): 494--509.

\bibitem[\citeproctext]{ref-cohen88}
Cohen, Jacob. 1988. \emph{Statistical Power Analysis for the Behavioral
Sciences}. 2nd ed. Lawrence Erlbaum.

\bibitem[\citeproctext]{ref-gabriel96}
Gabriel, Richard P. 1996. \emph{Patterns of Software: Tales from the
Software Community}. Oxford University Press.

\bibitem[\citeproctext]{ref-glass02}
Glass, Robert L. 2002. \emph{Facts and Fallacies of Software
Engineering}. Addison-Wesley.

\bibitem[\citeproctext]{ref-hess04}
Hess, Melinda R., and Jeffrey D. Kromrey. 2004. {``Robust Confidence
Intervals for Effect Sizes.''} \emph{Annual Meeting of the American
Educational Research Association}.

\bibitem[\citeproctext]{ref-hunt99}
Hunt, Andrew, and David Thomas. 1999. \emph{The Pragmatic Programmer}.
Addison-Wesley.

\bibitem[\citeproctext]{ref-knuth84}
Knuth, Donald E. 1984. {``Literate Programming.''} \emph{The Computer
Journal} 27 (2): 97--111.

\bibitem[\citeproctext]{ref-lions77}
Lions, John. 1996. \emph{Lions' Commentary on {UNIX} 6th Edition, with
Source Code}. Peer-to-Peer Communications.

\bibitem[\citeproctext]{ref-massey51}
Massey, Frank J. 1951. {``The {Kolmogorov-Smirnov} Test for Goodness of
Fit.''} \emph{Journal of the American Statistical Association} 46 (253):
68--78.

\bibitem[\citeproctext]{ref-menzies26ezr}
Menzies, Tim, and Srinath Srinivasan. 2026. {``Can {AI} Be Easy? Lessons
Learned from the {EZR.py} Toolkit.''} \emph{Software: Practice and
Experience}.

\bibitem[\citeproctext]{ref-norvig92}
Norvig, Peter. 1992. \emph{Paradigms of Artificial Intelligence
Programming: Case Studies in Common Lisp}. Morgan Kaufmann.

\bibitem[\citeproctext]{ref-ousterhout18}
Ousterhout, John. 2018. \emph{A Philosophy of Software Design}. Yaknyam
Press.

\bibitem[\citeproctext]{ref-park88}
Park, Stephen K., and Keith W. Miller. 1988. {``Random Number
Generators: Good Ones Are Hard to Find.''} \emph{Communications of the
ACM} 31 (10): 1192--201.

\bibitem[\citeproctext]{ref-pearl18}
Pearl, Judea, and Dana Mackenzie. 2018. \emph{The Book of Why: The New
Science of Cause and Effect}. Basic Books.

\bibitem[\citeproctext]{ref-schreiber00}
Schreiber, Guus, Hans Akkermans, Anjo Anjewierden, et al. 2000.
\emph{Knowledge Engineering and Management: The {CommonKADS}
Methodology}. MIT Press.

\bibitem[\citeproctext]{ref-seibel05}
Seibel, Peter. 2005. \emph{Practical Common Lisp}. Apress.

\bibitem[\citeproctext]{ref-spinellis03}
Spinellis, Diomidis. 2003. \emph{Code Reading: The Open Source
Perspective}. Addison-Wesley.

\bibitem[\citeproctext]{ref-stevens46}
Stevens, S. S. 1946. {``On the Theory of Scales of Measurement.''}
\emph{Science} 103 (2684): 677--80.

\bibitem[\citeproctext]{ref-welford62}
Welford, B. P. 1962. {``Note on a Method for Calculating Corrected Sums
of Squares and Products.''} \emph{Technometrics} 4 (3): 419--20.

\bibitem[\citeproctext]{ref-wilson22}
Wilson, Greg. 2024. \emph{Software Design by Example: A Tool-Based
Introduction with Python}. CRC Press.

\end{CSLReferences}

\backmatter
\end{document}
